\subsection{Течение жидкости}
%%%%%%%%%%%%%%%%%%%%%%%%%%%%%%%%%%%%%%%%%%%%%%%%%%%%
%%%%%%%%%%%%%%%%%%%%%%%%%%%%%%%%%%%%%%%%%%%%%%%%%%%%
 \z{Насос должен подавать ежесекундно объём воды $V$ на высоту $h$ по трубе постоянного сечения $S$. Какова должна быть мощность $P$ насоса? Плотность воды $\rho$.}
%
{$N = \rho V \left(gh + \dfrac{V^2}{2S^2}\right)$}
%%%%%%%%%%%%%%%%%%%%%%%%%%%%%%%%%%%%%%%%%%%%%%%%%%%%
\z{Насос представляет собой расположенный горизонтально цилиндр с поршнем площади $S$ и выходным отверстием площади $s$, расположенным на оси цилиндра. Определите скорость $v$ истечения струи жидкости из насоса, если поршень под действием силы $F$ перемещается с постоянной скоростью. Плотность жидкости $\rho$.}
%
{$v = \sqrt{\dfrac{2FS}{\rho(S^2 - s^2)}}$}
%%%%%%%%%%%%%%%%%%%%%%%%%%%%%%%%%%%%%%%%%%%%%%%%%%%% 
 \z{Плита массы $m$ удерживается на месте в горизонтальном положении $N$ струями жидкости плотности $\rho$, бьющими вертикально вверх. Площадь каждого отверстия $S$. Скорость жидкости на выходе из отверстий $v$. На какой высоте $H$ над отверстиями удерживается плита, если, достигнув плиты, жидкость разлетается от неё в горизонтальной плоскости?}
%
{$h = \dfrac{1}{2g} \left(v^2 - \left(\dfrac{mg}{\rho v S N}\right)^2\right)$}
%%%%%%%%%%%%%%%%%%%%%%%%%%%%%%%%%%%%%%%%%%%%%%%%%%%% 
 \z[{\begin{wrapfigure}[8]{r}{0.2\textwidth}%
		\vspace{-20pt}%
		\begin{flushright}%
 			\includegraphics[width=0.2\textwidth]{9.5.4}%
		\end{flushright}%
	\end{wrapfigure}%
	}]%
	{С каким ускорением будет двигаться длинное цилиндрическое тело плотности $\rho$ и радиуса $r$ вдоль оси вертикального высокого цилиндрического сосуда радиуса $R$, заполненного жидкостью плотности $\rho_0$? Чему равна разность давлений $\Delta p$ на верхнее и нижнее основания тела, если его длина равна $h$?}
%
{$a = \dfrac{\rho - \rho_0}{\rho + \dfrac{\rho_0 r^2}{R^2 - r^2}} g$; ~ $\Delta P = \dfrac{\rho R^2}{\rho(R^2 - r^2) + \rho_0 r^2} \rho_0 g h$.}
%%%%%%%%%%%%%%%%%%%%%%%%%%%%%%%%%%%%%%%%%%%%%%%%%%%%	
 \z {Из отверстия в дне высокого сосуда вытекает вода. Сечение сосуда $S$, сечение струи $s$. Уровень воды в сосуде перемещается с постоянным ускорением. Найдите это ускорение.}
%
{$a = g \left(\dfrac{s}{S}\right)^2$}